\chapter[Resultados: clasificación y detección de melanoma]{Resultados: clasificación diagnóstica y detección de melanoma}
\label{cp:res_clasif}
% Extensión estimada: 20--25 páginas.
% Fuentes: sección "Clasificación diagnóstica" del documento de experimentación, con los números del split
% disjunto por lesión (Codigo/outputs/results/lesion_split_report.json, cerrado el 2026-09-15 con las
% semillas 42/43/44).
%
% Feedback del tutor (14/09/2026): la sección de cierre se llama "Discusión" en lugar de "Síntesis del
% capítulo". La discusión transversal, que cruza este capítulo con el de generación, queda en el
% Capítulo~\ref{cp:discusion}.

\pendiente{Este capítulo responde la PI1 y contrasta las hipótesis H1, H3 y H4. Los resultados se presentan
en el orden en que se tomaron las decisiones: cada sección presenta un resultado, lo interpreta y justifica
el paso siguiente.}

\section{Clasificadores supervisados}
\label{sec:res_supervisados}
\pendiente{Se presentan el barrido y el modelo elegido en validación, sus resultados en test como media
$\pm$ desvío entre semillas (exactitud, F1 ponderado y macro, sensibilidad macro y \emph{recall} por clase),
el efecto de las estrategias de desbalance y el efecto de corregir la fuga respecto del split original. Se
discute también la dispersión entre semillas, muy desigual según la métrica: la exactitud varía poco mientras
que el \emph{recall} de melanoma es la métrica más inestable, lo que obliga a reportar media y desvío y no
una corrida única.}

\section{Modelos visión--lenguaje médicos generales}
\label{sec:res_vlm_generales}
\pendiente{Se comparan las cuatro estrategias con Mini-VLM y LLaVA-Med frente al baseline supervisado, y se
muestra el colapso de las predicciones hacia una o dos clases.}

\section{Re-diagnóstico del colapso}
\label{sec:res_rediagnostico}
\pendiente{Se contrasta la hipótesis del desbalance con la evidencia que la refuta: el colapso aparece en
\emph{zero-shot}, los datos de ajuste no están sesgados hacia las clases problemáticas y cada configuración
colapsa hacia clases distintas.}

\section{La especialización de dominio recupera la discriminación}
\label{sec:res_dermlip}
\pendiente{Se presentan los resultados de DermLIP en \emph{zero-shot}, \emph{linear probe}, kNN y SVM (split
por lesión) frente a los VLMs generales y a los clasificadores supervisados.}

\section{Comparación con el estado del arte}
\label{sec:res_sota_clasif}
\pendiente{Se comparan los resultados con los publicados sobre HAM10000 (Sección~\ref{sec:estado_arte}),
señalando las diferencias de protocolo (particiones, fuga, métricas) que limitan la comparación.}

\section{Calibración y significancia estadística}
\label{sec:res_calib_stats}
\pendiente{Se analiza la calibración antes y después de \emph{temperature} y \emph{vector scaling}, junto con
el efecto de los pesos por clase, y se reportan los intervalos de confianza y las pruebas estadísticas entre
los modelos finales.}

\section{Robustez}
\label{sec:res_robustez}
\pendiente{Se evalúa la generalización a HIBA (otra población y otros fototipos) y a perturbaciones de las
condiciones de adquisición.}

\section{Interpretabilidad}
\label{sec:res_interpretabilidad}
\pendiente{Se analizan mapas Grad-CAM para determinar si el modelo atiende a la lesión.}

\section{Detección de melanoma}
\label{sec:res_melanoma}
\pendiente{Se evalúan las tareas binarias melanoma-vs-resto y maligno-vs-benigno (AUC, sensibilidad y
especificidad con umbral fijado en validación) por modelo y para el ensamble, y se comparan con referencias
clínicas y de la literatura.}

\section{Discusión}
\label{sec:res_clasif_discusion}
\pendiente{Se resume qué hipótesis se sostienen y cuáles no, y por qué el clasificador especializado pasa a
ser la fuente de diagnóstico del sistema de generación.}
